\chapter{La Biblioteca \\texttt{python-template.lp}}
\label{chap:python-template}
{\huge$\bigstar$}{\huge$\bigstar$}{\huge$\bigstar$} Avanzado

\section{Visión General}
\begin{lispcode}
(load "lib/python-template.lp")
(write-python-template `(def ,'fib (n) "return fib(n-1) + fib(n-2)"))
\end{lispcode}

\section{Mapeo de Constructos}
\begin{tabular}{ll}
\toprule
Forma Lisp & Código Python \\
\midrule
\texttt{(def name (args) body)} & \texttt{def name(args):\textbackslash\textbackslash indent body} \\
\texttt{(return expr)} & \texttt{return expr} \\
\texttt{(if c t e)} & \texttt{if c:\textbackslash\textbackslash indent t\textbackslash\textbackslash else:\textbackslash\textbackslash indent e} \\
\texttt{(for v in e body)} & \texttt{for v in e:\textbackslash\textbackslash indent body} \\
\texttt{(:= v e)} & \texttt{v = e} \\
\texttt{(+ a b)} & \texttt{a + b} \\
\bottomrule
\end{tabular}

\begin{examplebox}[title=Generar función factorial]
\begin{lispcode}
(write-python-template
  `(def ,'factorial (n)
     "result = 1"
     (for i in (range 1 (+ n 1))
       (:= result (* result i)))
     (return result)))
\end{lispcode}
\end{examplebox}

